%
% Documento: Resumo (Português)
%

\section{\textbf{Resumo}}
  Esta pesquisa visa contribuir para a determinação de um parâmetro fundamental em astronomia: a massa de estrelas jovens em aglomerados estelares, com potencial para enriquecer a compreensão da rotação estelar. O objetivo é estimar massas estelares em aglomerados como o da Nebulosa de Órion (ONC), Plêiades, h Persei, Upper Scorpius e NGC 2516 utilizando uma combinação de métodos tradicionais, como ajuste de isócronas e relações massa-luminosidade, e técnicas de aprendizado de máquina. Adotamos uma abordagem descritiva, utilizando ferramentas estatísticas para avaliar a qualidade das estimativas. A amostra é não probabilística, com foco em estrelas jovens de baixa massa, e os dados foram coletados de diversos catálogos, incluindo GAIA e TESS. As magnitudes absolutas foram calculadas a partir da fotometria, e a extinção interestelar individual foi considerada para os aglomerados mais jovens. Os resultados foram comparados com valores de massa obtidos recentemente na literatura. As massas calculadas usando a relação massa-luminosidade apresentaram uma boa concordância com a maioria dos aglomerados, principalmente os com uma sequência principal bem definida, para as Plêiades, e do ajuste de isócronas para aglomerados com mais objetos na pré-sequencia principal, como a ONC, embora com incertezas e limitações inerentes. A integração de técnicas de aprendizado de máquina se mostrou promissora para aumentar a precisão e confiabilidade das estimativas de massa estelar, apontando direções futuras para pesquisas astrofísicas.
